\lysh{21447}{4}
\begin{alist}
\item Ο αριθμός των βακτηρίων όταν ξεκίνησε το πείραμα ήταν $P(0) = 200\cdot e^{c\cdot 0} = 200$ βακτήρια.
\item Έχουμε:
\[P(1) = 328 \Leftrightarrow 200\cdot e^{c\cdot 1} = 328 \Leftrightarrow e^c = \frac{328}{200} \Leftrightarrow e^c = 1{,}64 \Leftrightarrow c = \ln(1{,}64) \Leftrightarrow c = 0{,}5.\]
Άρα $c = \dfrac{1}{2}$.
\item Ο αριθμός των βακτηρίων είναι μεγαλύτερος από το δεκαπλάσιο και μικρότερος από το εκατονταπλάσιο της αρχικής του τιμής, δηλαδή
\begin{align*}
10\cdot P(0) &< P(t) < 100\cdot P(0) \Leftrightarrow \\
10\cdot 200 &< 200\cdot e^{\frac{1}{2}\cdot t} < 100\cdot 200 \Leftrightarrow \\
10 &< e^{\frac{1}{2}\cdot t} < 100 \overset{\ln x \nearrow}{\iff} \\
\ln 10 &< \ln(e^{\frac{1}{2}\cdot t}) < \ln 100 \Leftrightarrow \\
\ln 10 &< \frac{1}{2}\cdot t < \ln 10^2 \Leftrightarrow \\
2\cdot\ln 10 &< t < 4\cdot\ln 10 \Leftrightarrow \\
2\cdot 2{,}3 &< t < 4\cdot 2{,}3 \Leftrightarrow \\
4{,}6 &< t < 9{,}2.
\end{align*}
Άρα το ζητούμενο χρονικό διάστημα (σε ώρες) είναι $4{,}6 < t < 9{,}2$.
\end{alist}

